\section{Current Implementation and Evidence Boundary}

Resolution-aware search is a package contract as much as an algorithm choice.
The current evidence path must name the unit type, fine-to-coarse mapping,
adjacency derivation, population aggregation, and any information lost in the
chosen resolution.  A manifest field is useful only if a verifier can replay the
mapping and detect mismatched unit order or topology.

Claims about mixing improvement or legal adequacy should therefore be tied to a
specific archived package.  Resolution changes can improve computational
behavior, but they do not by themselves prove that the resulting plan is
legally sufficient or comparable to tract-level runs.

The checked-in \texttt{U.11+resolution-smoke} package is the minimal positive
case for this boundary.  It verifies GEOID-prefix mapping, manifest resolution
fields, population aggregation, and derived coarse adjacency in
\texttt{bisect-multiscale::resolution\_evidence}.  It deliberately stops short
of claiming the Texas autocorrelation reduction, BG-level compactness gains, or
court-ready legal adequacy.
